The boats went out at the turn of the tide and she was on the stones to see them go, the way she had been on the stones to see them go since she was nine years old, except that now they went when she said.

Nobody had decided that. There had been no moment. What there had been was the third morning after the fighting, when four boats were fit to swim and eleven men wanted to be in them and the argument on the Hook had gone on so long that they missed the tide entirely, and Wyla had stood in the middle of it with her arms folded and said: \textit{Corrin’s boat, Corrin and the Flint lads and Old Mott’s grandson, and the rest of you get tomorrow, and if you can’t hold to that then I’ll go and sit down and you can all starve arguing.}

They had held to it. She had not expected them to and they had, and the next morning eleven men had come and stood in front of her without being asked, and she had done it again.

That was the whole of it. That was the entire construction, top to bottom, and it had taken about ninety heartbeats on a wet quay and there was not one thing written down anywhere in the world about it.

Nine boats now. Six that swam properly and three that wanted bailing all day, and one of the three was Corrin’s cousin’s and had a hole in it stuffed with a shirt, and she let it out anyway because Corrin’s cousin had four in his house and no other way to feed them.

\textit{And because he’ll not forget it,} she thought, watching the shirt-boat go out low in the water. \textit{That’s the other half.}

\sceneBreak

The Hook had come through better than most of the city, which was a thing nobody said out loud in front of people from Flea Bottom.

No cellars. That was all it was. You could not have a cellar within two hundred feet of that water without it filling twice a day, so nobody on the Hook had ever dug one, so on the second night when the dead came up out of the ground all along the lower city they had not come up on the Hook, because there was nowhere for them to come up from. Half the fish sheds still had roofs. There were eleven houses standing on Rope Lane with people in them who had gone to bed on the first night and woken up on the fourth and lost nobody at all.

Wyla had lost Old Mott’s daughter and a boy called Pyat who used to gut for her and a woman she had known forty years whose name she was ashamed to find she had to think about.

Three. Three out of a lane. There were streets up the hill that had lost everybody.

She did not think about that too hard. It was not the sort of thing that went anywhere useful.

\sceneBreak

The market was up and running by the second hour and the whole of it was talking about the eunuchs.

She had heard it four times before she got to the top of the fish lane, in four different shapes, and the shapes were getting better as they went — that was the thing about a story, it got \textit{tidier}, you could hear the corners coming off it. By the fourth telling the eunuchs had been drawn up in a line watching, and by the time a woman from the Street of Sisters said it to her face at the top of the lane they had been drawn up in a line watching and one of them had said something in their own tongue and the rest had laughed.

“They never,” Wyla said.

“They did. Mag’s sister’s boy was on the eastern tier and he heard it.”

“Heard what? He speak Ghiscari, does he, Mag’s sister’s boy?”

The woman looked at her.

“He heard them \textit{laugh}, Wyla.”

“Aye, well.” Wyla set down the basket. “I’ve heard my Corrin laugh at a funeral and it weren’t because he found it funny. Man laughs when he’s frighted. So do women. So do you — you laughed at me the night the bells went and I’ve never once thought you were enjoying yourself.”

The woman went off with her mouth in a line, and Wyla thought about it for as long as it took to pick the basket up, which was about the right amount of time, and then let it go.

\textit{It’ll not matter what’s true,} she thought. \textit{That story’s got legs on it now. It’s in a shape people like.}

And then she stopped thinking about it, because Corrin’s boat was in early with a good catch and there were four things she needed to be doing before the second bell.

\sceneBreak

She got the whole of the market’s business by standing still and letting it come to her, which was a trick she had learned from her mother and which she had never once had to explain to anybody because nobody had ever asked.

By noon she knew: that the pits had opened again at first light and gone quiet and orderly and that there had been a wagon put crossways at the top of the ramp, which everybody agreed was sensible and nobody could say who had thought of. That the big woman knight who had died in the arch had been put in the first load with her name read out, in the common pit, same as the rest — which Wyla heard three times and which nobody in that market found remarkable in the slightest, though a fishwife did say it was a waste of a good sword and got told the sword had gone home on a ship. That two of the wells on the Street of Steel had gone brackish. That a man on Cobbler’s Square was selling boiled water at a copper a cup and had a queue.

That last one she went and looked at, because she wanted to see it.

He had a fire, which was the whole of it. Somebody had got hold of wood — she did not ask, because she knew what asking got you — and he was boiling bay water down in a great iron pan and selling what came off the lid, and it was a filthy slow way of doing it and it worked, and there were forty people in his queue with cups.

\textit{Copper a cup,} she thought. \textit{That’s robbery.}

Then she stood and watched it a while longer and revised.

\textit{No. That’s a man with a fire and one pan doing what one pan does, and if he charged less he’d be out of wood by tomorrow and there’d be nobody boiling anything. That’s not robbery, that’s a copper a cup.}

She went and found out where he got the wood.

\sceneBreak

The water tally was being cheated on the Hook and had been since about the sixth day, and Wyla knew exactly who was doing it.

It was the Flints. It was three houses of Flints who had put themselves down as four houses of Flints, because the boy taking the count on the second day had been about fifteen and frightened and had written what he was told, and the Flints had been drawing a fourth house’s water for a week and selling the difference up the hill.

She had known for four days and had said nothing to anybody.

She thought about it that afternoon while she was doing the gutting, which was when she thought about most things.

She could go to the lord’s clerk. That was one way. It would take an hour and the Flints would lose the fourth house and there would be a good deal more water on the Hook, and every Flint on that lane would know inside a day who had done it, and there were nine of them and two were on boats.

Or.

She got Elder Flint on his own on the quay at dusk and did not raise her voice at all.

“Fourth house,” she said.

He went the colour of a fish belly.

“Now don’t.” She held a hand up. “I’ve not gone to anybody and I’m not going. I want two things and then it’s finished and I’ll not say it again.” She counted them on her fingers, which she did when she wanted a thing remembered. “One: the fourth house’s water comes to me from tomorrow, every day, and I’ll say where it goes. Two: your Rennick and your young Tom go out in the shirt-boat with Corrin’s cousin for a fortnight and they bail, and they don’t complain about it, and if I hear one word out of either of them I’ll go to the clerk before the words are cold.”

Elder Flint stood there working it out.

“That’s not nothing you’re asking.”

“It’s a good deal less than the clerk.”

“Aye,” he said, after a moment. “It is that.”

She watched him go, and thought: \textit{That’s better all round. That’s more water where it’s wanted and two lads on a boat that needs them and nobody hanged over a fourth house.}

It did not once occur to her, then or later, that she had just taken over the water on the Hook.

\sceneBreak

There was one part of the day she did not much like and it came at about the fourth hour, when the boats were in and the baskets were on the stones and there were more people than baskets.

There were always more people than baskets. The whole city knew now that the Hook had fish and the whole city had begun walking down the hill for it, and by the fourth hour there might be sixty or eighty people standing on those stones with nothing to trade but a promise, and she had perhaps forty baskets.

She did it herself. She had tried letting Corrin do it for two days and he had given it all away in an hour and then stood there with his hands out like a man who had been robbed.

So it was her, and she went along the line, and she went fast, because going slow was worse for everybody.

\textit{Sara Wick’s — yes, she’s got the three little ones and she carried for me on the fourth night. Malliver — yes, he mends nets and I’ll want nets. That one, no, I don’t know her. That one, no, and don’t tell me about your mother, everybody here’s got a mother. Him, yes, he’s Corrin’s wife’s brother and I’ll hear about it for a month if I don’t. Her, no. Her — }

She stopped at a girl of about seventeen holding a baby, and did not know her, and had never seen her on the Hook in her life.

“Where are you from?”

“Pisswater Bend.”

“Then why aren’t you at Pisswater Bend?”

“There isn’t one,” the girl said.

Wyla gave her a basket.

She went on down the line, and gave out the rest, and turned twenty-odd people away, and one of them was a man of about fifty who did not argue at all and simply stood there for a while afterward looking at the stones before he went back up the hill.

She would have said, if anybody had asked her, that she gave it out by need. That is what she would have said and she would have said it honestly, and she would have been able to name for you the need of every single person she gave a basket to that afternoon, which was true of every one of them, and there was not one word of a lie in any of it.

Nobody asked her.

\sceneBreak

Lord Davos came down on the eleventh day and she made him wait.

Not out of spite. She had four boats to get in and a dispute between two of Old Mott’s grandsons about a net, and she saw him come along the quay in that coat of his and she thought \textit{he’ll keep} and got on with the net.

It was about twenty minutes. He stood at the end of the stones the whole time with his hands behind his back and did not once try to catch her eye, and when she finally came over he said, “That’s a good net,” which she thought was the correct thing to have said and probably not calculated, which was better.

“What do you want, my lord.”

“Fish.”

“You and the rest of the city.”

“I want it in a different way from the rest of the city.” He got a paper out and then put it back, which she noticed. “The grain’s finished by the turn of the moon. That’s not a guess, that’s the western stores counted twice. There’s a ship coming up from the Reach if the Reach still has anything to send, and there’s four going south, and none of it lands before the moon turns. Between now and then this city eats what comes out of that water or it doesn’t eat.”

“Then you’d best hope the weather holds.”

“I’d like to put eleven more boats in it.”

Wyla laughed at him. It came out before she could think about whether it was wise, and she found she did not much mind.

“Where from?”

“There’s hulls in the Blackwater that’d swim with work. There’s Greyjoy ships that’d tow. There’s men in this city with nothing to do, and — “

“My lord.” She wiped her hands on the front of her. “You’ve got hulls. Grand. Who’s crewing them?”

Davos Seaworth stopped.

“There’s it,” she said. “You can find me eleven hulls by Tuesday and I’ll not have a man to put in one of them, because a boat isn’t a barrel, it’s four men who know each other, and there’s nine boats going out of this Hook and there’d be nine going out if you gave me fifty hulls, because there’s thirty-six men left on these stones who can work a boat and I know every single one of their names and I know which of them can stand each other.” She let that sit. “There’s your fleet, my lord. It’s thirty-six men and it’s not getting bigger.”

Davos was quiet for a moment.

“Could you teach them?”

“Teach who?”

“Anyone. Boys. Women. There’s three thousand people sitting in the Dragonpit doing nothing but being frightened.” He said it slowly, the way a man says a thing he is working out while he says it. “You’d not want them on a boat. But you’d want them gutting, and salting, and mending, and hauling, and carrying it up the hill, and if all that came off your thirty-six men then your thirty-six men would be out on the water an extra two hours a day, and that’s — “

He stopped, because she had started nodding about halfway through, and he had noticed.

“That’s the whole of it,” Wyla said. “That’s what I’ve been trying to say to Corrin for a week and he looks at me like a cod.” She was already turning it over. “I’d want the salt. There’s no salt. You can catch what you like and it rots by Thursday, and there’s no salt in this city since the Reach road shut and I’ve been drying on the racks like it’s a hundred years ago.”

“How much salt?”

“How much can you get?”

“I don’t know yet.”

“Then find out and come back,” Wyla said, “and don’t bring me a paper.”

\sceneBreak

He came back four hours later with an answer about the salt, which was that there was some, in the wrong place, held by somebody difficult, and that he would have it in six days or not at all.

And then he stood on the stones a bit longer, and she understood that there was a second thing.

“Say it, then.”

“The barges,” Davos said. “They start on the fourteenth day. Whatever the pits can’t take goes out on Greyjoy hulls to deep water past the point.”

“Aye. I’d heard.”

“You’d heard.”

“My lord, I’ve nine boats on that water every morning. Course I’ve heard.” Wyla looked out past the mole, where the light was going, and thought about how to say it, and then said it the way she said everything. “You want to know if the fish’ll be fit.”

“Yes.”

“They’ll be fit. It’s past the point and it’s deep and it’s cold and that water doesn’t come back in on this tide, it goes east. My grandfather’s grandfather fished that ground and there’s been ships going down out there since there’s been ships.” She shrugged. “It’s the sea. There’s things in it. There’s always been things in it.”

Davos nodded slowly.

“And will people eat it?”

And there it was, which was the actual question, and she liked him better for having got to it.

“Not if you tell them,” she said.

He did not answer.

“Don’t look at me like that. I’m not saying lie to them, I’d not say that, and anyway there’s no lying to a market — there’s a hundred people on that quay every morning and half of them can smell a thing before you’ve said it.” She picked up the basket. “I’m saying don’t \textit{announce} it. Don’t have a man stand up in the Dragonpit and say \textit{the fish is safe}, because the day you say the fish is safe is the day nobody in this city eats fish again as long as they live. Send the barges out past the point, and don’t put it in a proclamation, and in a fortnight there’ll not be a soul who thinks about it.”

“And in a year?”

“In a year,” said Wyla, “there’ll be a story about it, and it’ll be a good one, and I’ll not be able to sell a mackerel off that ground for love nor money.” She hitched the basket up onto her hip. “But we’ll all have got to the year, my lord, and I’d sooner have the story than the alternative.”

\sceneBreak

She slept well that night, the way she had slept every night since the second one.

That was the thing she never mentioned to anybody, not because she was ashamed of it but because there was nobody it made sense to say it to. For ten years she had lain down every night of her life and gone through it — the tower, the stairs, the asking, the man’s face when he said he had no memory of the name — and had got up in the morning with it still going, like a mill nobody had switched off.

It had stopped. On the third night she had noticed it had stopped, and had lain there in the dark waiting for it to start again, and it had not, and it had not the next night either.

She did not turn it over. She did not think about the kitchens, or about who had thrown the first thing, or about the noise, or about whether the man had saved forty people that night, which he probably had. She had been told once by a septa that a person who does not examine their soul will find it examined for them, and she had been nine years old at the time and had thought it sounded like something a woman said who had never had to sell fish.

What she thought about, lying on her back on Rope Lane on the eleventh night, was salt.

Six days for the salt, and the barges on the fourteenth, and eleven hulls she could not crew, and three thousand people in the Dragonpit who could gut if somebody showed them, and the Flints’ fourth house of water, and whether Sara Wick’s eldest was strong enough for the racks yet, and whether the man on Cobbler’s Square would sell her the name of whoever had the wood.

She was asleep before she got to the wood.